\BourbakiExerciseGroup

\begin{enumerate}
\def\labelenumi{\arabic{enumi})}
\tightlist
\item
  \begin{enumerate}
  \def\labelenumii{\alph{enumii})}
  \tightlist
  \item
    设 \(q\) 是某个素数的幂。证明：对于每个整数 \(n > 1\) ，若 \(h_i\) （ \(1 \leqslant i \leqslant r\) ）是整数 \(n\) 的全部不同素因子，则满足 \(\zeta \in \mathbf{F}_{q^n}\) 且 \(\mathbf{F}_{q^n} = \mathbf{F}_q(\zeta)\) 的元素的个数等于
  \end{enumerate}
\end{enumerate}

\[
\nu = q ^ {n} - \sum_ {i} q ^ {n / h _ {i}} + \sum_ {i <   j} q ^ {n / h _ {i} h _ {j}} - \sum_ {i <   j <   k} q ^ {n / h _ {i} h _ {j} h _ {k}} + \dots + (- 1) ^ {r} q ^ {n / h _ {1} h _ {2} \dots h _ {r}}
\]

（注意，这样的元素由不属于任何一个域 \(F_{q}^{n/h_{i}}\) 的性质来刻画）。若 \(h_{1} \leqslant \cdots \leqslant h_{r}\) ，我们有

\[
q ^ {n} - \sum_ {i} q ^ {n / h _ {i}} \leqslant \nu \leqslant q ^ {n} - q ^ {n / h _ {1}}.
\]

考察 \(n\) 是素数之幂的情形。

\begin{enumerate}
\def\labelenumi{\alph{enumi})}
\setcounter{enumi}{1}
\item
  设 \(b_{n}(q)\) 为次数为 \(n\) 、位于 \(\mathbf{F}_q\) 上的首一不可约多项式的个数。试由 \(a\) ) 导出 \(b_{n}(q)\) 的一个表达式。特别地，若 \(q \geqslant 3\) ，证明次数为 \(n\) 、属于 \(\mathbf{F}_q[\mathbf{X}]\) 的不可约多项式（未必首一）的个数至少等于 \(q^n (q - 2) / n \geqslant q^{n + 1} / 3n\) 。
\item
  证明 \(\sum_{d\mid n}db_d(q) = q^n\) （建立恒等式 \(1 / (1 - q\mathrm{T}) = \prod_{n}(1 - \mathrm{T}^n)^{-b_n(q)}\) ）。
\end{enumerate}

* 利用默比乌斯反演公式（Lie, II, 第 176 页，附录），证明

\[
n b _ {n} (q) = \sum_ {d \mid n} \mu (n / d) q ^ {d}.
\]

将这一结果与 \(b\) 中所得的结果比较。 *

\begin{enumerate}
\def\labelenumi{\alph{enumi})}
\setcounter{enumi}{3}
\tightlist
\item
  设 \(a_{n}(q)\) 为 \(\mathbf{F}_q[\mathbf{X}]\) 中无重因式的首一多项式的个数。证明 \(a_0(q) = 1, \quad a_1(q) = q, \quad a_n(q) = q^n - q^{n-1}\) 对 \(n > 1\) 成立（建立恒等式 \(q^n = \sum_i a_{n-2i}(q) q^i\) ）。
\end{enumerate}

\begin{enumerate}
\def\labelenumi{\arabic{enumi})}
\setcounter{enumi}{1}
\item
  证明次数为 \(n\) 、属于 \(\mathbf{F}_q[\mathbf{X}]\) 且在 \(\mathbf{F}_q\) 中没有根的首一多项式的个数是 \(q^{n - q}(q - 1)^q\) ，其中 \(n \geqslant q\) （注意每个函数 \(\mathbf{F}_q \to \mathbf{F}_q\) 都是多项式函数）。
\item
  设 \(\mathbf{K}\) 是含 \(q\) 个元素的有限域；对每个素数 \(l\) ，令 \(N_{l}\) 为 \(\mathbf{K}\) 的满足下列条件的扩张的并：它们包含在 \(\Omega\) —— \(\mathbf{K}\) 的一个代数闭包 —— 之中，且次数均为 \(l\) 的幂。证明 \(\mathbf{N}_{l}\) 是 \(\mathbf{K}\) 的一个阿贝尔扩张，其伽罗瓦群同构于群 \(\mathbf{Z}_{l}\) ，即 \(l\) 进整数所成的群。由此推出 \(\Omega\) 同构于张量积 \(\otimes \mathbf{N}_{l}\) ，从而
\end{enumerate}

\(\operatorname{Gal}(\Omega / \mathbf{K})\) 同构于乘积 \(\prod_{l} \mathbf{Z}_l\) 。

* 4) 将多项式 \(\mathbf{X}^4 + 1 \in \mathbf{F}_p[\mathbf{X}]\) 分解为不可约因式。 *

\begin{enumerate}
\def\labelenumi{\arabic{enumi})}
\setcounter{enumi}{4}
\tightlist
\item
  设 \(\mathbf{K}\) 是含 \(q\) 个元素的有限域，其中 \(q\) 不是 2 的幂，并设 \(a_1, a_2, b\) 是 \(\mathbf{K}\) 的三个元素，使得 \(a_1a_2 \neq 0\) 。证明：数目 \(\nu\) （即解 \((x_1, x_2) \in \mathbf{K}^2\) 的数目，这些解满足方程 \(a_1x_1^2 + a_2x_2^2 = b\) ）由下列公式给出：
\end{enumerate}

\(1^{\circ}\) 若 \(b = 0\) 且 \(-a_{1}a_{2}\) 在 \(\mathbf{K},\nu = 1\)

\(2^{\circ}\) 若 \(b \neq 0\) 且 \(-a_{1}a_{2}\) 在 \(\mathbf{K}\) 中不是平方元，则 \(\nu = q + 1\) ；

\(3^{\circ}\) 若 \(b = 0\) 且 \(-a_{1}a_{2}\) 是 \(\mathbf{K}\) 中的平方元，则 \(\nu = 2q - 1\) ；

\(4^{\circ}\) 若 \(b \neq 0\) 且 \(-a_{1}a_{2}\) 是 K 中的平方元，则 \(\nu = q - 1\) 。

（当 \(-a_{1}a_{2}\) 不是平方元时，向 K 添加 \(\mathbf{X}^{2} + a_{1}a_{2}\) 的一个根，并利用有限域的乘法群是循环群这一事实。）

\begin{enumerate}
\def\labelenumi{\arabic{enumi})}
\setcounter{enumi}{5}
\tightlist
\item
  设 \(\mathbf{K}\) 是含 \(q\) 个元素的有限域，其中 \(q\) 不是 2 的幂。
\end{enumerate}

\begin{enumerate}
\def\labelenumi{\alph{enumi})}
\tightlist
\item
  设 \(a_{1}, a_{2}, \ldots, a_{2m}, b\) 是 K 的元素，使得 \(a_{1}a_{2} \ldots a_{2m} \neq 0\) 。证明：数目 \(\nu\) （即解 \((x_{1}, \ldots, x_{2m}) \in \mathbf{K}^{2m}\) 的数目，这些解满足方程 \(a_{1}x_{1}^{2} + \cdots + a_{2m}x_{2m}^{2} = b\) ）由下列公式给出：
\end{enumerate}

\(1^{\circ}\) 若 \(b = 0\) 且 \((-1)^{m}a_{1}a_{2}\ldots a_{2m}\) 在 K 中不是平方元，则 \(\nu = q^{2m - 1} - q^{m} + q^{m - 1}\) ；

\(2^{\circ}\) 若 \(b \neq 0\) 且 \((-1)^{m}a_{1}a_{2} \ldots a_{2m}\) 在 K 中不是平方元，则 \(\nu = q^{2m-1} + q^{m-1}\) ；

\(3^{\circ}\) 若 \(b = 0\) 且 \((-1)^{m}a_{1}a_{2}\ldots a_{2m}\) 是 K 中的平方元，则 \(\nu = q^{2m - 1} + q^{m} - q^{m - 1}\) ；

\(4^{\circ}\) 若 \(b \neq 0\) 且 \((-1)^{m}a_{1}a_{2} \ldots a_{2m}\) 是 K 中的平方元，则 \(\nu = q^{2m - 1} - q^{m - 1}\) 。

（对 \(m\) 作归纳论证，并利用习题 5。）

\begin{enumerate}
\def\labelenumi{\alph{enumi})}
\setcounter{enumi}{1}
\tightlist
\item
  设 \(a_1, a_2, \ldots, a_{2m+1}\) ，\(b\) 为 \(K\) 的元素，使得 \(a_1a_2 \ldots a_{2m+1} \neq 0\) 。证明：数目 \(\nu\) ，即下方方程的解 \((x_1, x_2, \ldots, x_{2m+1}) \in K^{2m+1}\) 的数目，
\end{enumerate}

\[
a _ {1} x _ {1} ^ {2} + \dots + a _ {2 m + 1} x _ {2 m + 1} ^ {2} = b
\]

由以下公式给出：

\(1^{\circ}\) 若 \(b = 0, \nu = q^{2m}\) ；

\(2^{\circ}\) 若 \(b \neq 0\) 且 \((-1)^{m}a_{1}a_{2} \ldots a_{2m+1}b\) 在 K 中不是平方元，则 \(\nu = q^{2m} - q^{m}\) ；

\(3^{\circ}\) 若 \(b \neq 0\) 且 \((-1)^{m}a_{1}a_{2} \ldots a_{2m+1}b\) 是 \(\mathbf{K}\) 中的平方元，则 \(\nu = q^{2m} + q^{m}\) 。（归结为 \(a\) ）的情形。）

\begin{enumerate}
\def\labelenumi{\arabic{enumi})}
\setcounter{enumi}{6}
\tightlist
\item
  设 K 为有限域，E 为 K 的 n 次代数扩张，且 \((\alpha_{1}, \ldots, \alpha_{n})\) 为 E 在 K 上的正规基；假设循环伽罗瓦群 \(\operatorname{Gal}(\operatorname{E}/\operatorname{K})\) 由循环置换 \(\sigma = (\alpha_{1}\alpha_{2} \ldots \alpha_{n})\) 生成。设 \(\tau\) 为向量 K-空间 E 的一个自同构，使得对每个 \(x \in E\) ，\(\tau(x)\) 在 E 上与 x 共轭（参见 V，第 154 页，习题 8）。
\end{enumerate}

\begin{enumerate}
\def\labelenumi{\alph{enumi})}
\item
  证明：若 \(\mathbf{K}\) 至少有 3 个元素，则 \(\tau\) 是域 E 的 K-自同构（归结到 \(\tau(\alpha_1) = \alpha_1\) 的情形，并考虑 \(\tau(\alpha_1 + c\alpha_j)\) ，其中 \(c \in \mathbf{K}\) ）。
\item
  如果 \(K = F_{2}\) ，证明对每个元素 \(\rho \in \operatorname{Gal}(E/K)\) ，关系式 \(\rho(\alpha_{r}) = \alpha_{r'}\) 与 \(\rho(\alpha_{s}) = \alpha_{s'}\) 蕴含 \(r' - r \equiv s' - s \pmod{n}\) 。由此推出，若 n \textgreater{} 5，则 \(\tau\) 是域 E 的 K-自同构。（注意：若 \(\tau(\alpha_{1}) = \alpha_{1}\) 且 \(\tau(\alpha_{1} + \alpha_{h}) = \alpha_{1} + \alpha_{k}\) ，其中 \(k \neq h\) ，则必有 \(h + k = n + 2\) ；由此推出必然有 \(\tau(\alpha_{2}) = \alpha_{n}\) 且 \(\tau(\alpha_{4}) = \alpha_{4}\) 或 \(\tau(\alpha_{4}) = \alpha_{n-2}\) ，这导致矛盾。）
\item
  对于 \(\mathbf{K} = \mathbf{F}_2\) 且 \(3 \leqslant n \leqslant 5\) ，向量 K-空间 E 的下列自同构不是 E 的域自同构：
\end{enumerate}

\[
\text { for } \quad n = 3, \tau : (\alpha_ {1}, \alpha_ {2}, \alpha_ {3}) \mapsto (\alpha_ {1}, \alpha_ {3}, \alpha_ {2})
\]

\[
\text { for } \quad n = 4, \tau : (\alpha_ {1}, \alpha_ {2}, \alpha_ {3}, \alpha_ {4}) \mapsto (\alpha_ {1}, \alpha_ {4}, \alpha_ {3}, \alpha_ {2})
\]

\[
\text { for } \quad n = 5, \tau : (\alpha_ {1}, \alpha_ {2}, \alpha_ {3}, \alpha_ {4}, \alpha_ {5}) \mapsto (\alpha_ {1}, \alpha_ {5}, \alpha_ {4}, \alpha_ {3}, \alpha_ {2})
\]

但它们满足 \(\tau(x)\) 与 \(x\) 共轭，对一切 \(x \in E\) 成立。

\begin{enumerate}
\def\labelenumi{\arabic{enumi})}
\setcounter{enumi}{7}
\tightlist
\item
  设 \(\mathbf{K}\) 是含 \(q\) 个元素的有限域，并且
\end{enumerate}

\[
\mathrm{P} (\mathbf {X}) = a _ {0} + a _ {1} \mathbf {X} + \dots + a _ {q - 2} \mathbf {X} ^ {q - 2}
\]

K{[}X{]} 中次数为 q - 2 的一个多项式。设 r 为矩阵

\[
A = \left( \begin{array}{c c c c} a _ {0} & a _ {1} & \dots & a _ {q - 2} \\ a _ {1} & a _ {2} & \dots & a _ {0} \\ \dots & \dots & \dots & \dots \\ a _ {q - 2} & a _ {0} & \dots & a _ {q - 3} \end{array} \right).
\]

证明 P 在 K 中的、异于 0 的不同根的个数是 q - 1 - r。（考虑乘积 VA，其中 V 是元素 \(x_{1}, x_{2}, \ldots, x_{q-1}\) （属于 \(K^{*}\) ）的范德蒙德矩阵。）

\begin{enumerate}
\def\labelenumi{\arabic{enumi})}
\setcounter{enumi}{8}
\tightlist
\item
  设 \(\mathbf{K}\) 是含 \(q\) 个元素的有限域，并令 \(r = q^n\) ，其中 \(n \geqslant 1\) 为整数。对每个多项式 \(\mathrm{P}(\mathbf{X}) = \sum_{j=0}^{m} a_j \mathbf{X}^j\) （属于 \(\mathbf{K}[\mathbf{X}]\) ），令
\end{enumerate}

\[
\hat {\mathrm{P}} (\mathrm{X}) = \sum_ {j} a _ {j} \mathrm{X} ^ {(r ^ {j} - 1) / (r - 1)}
\]

和

\[
\tilde {\mathrm{P}} (\mathrm{X}) = \mathrm{X} \hat {\mathrm{P}} \left(\mathrm{X} ^ {r - 1}\right) = \sum_ {j} a _ {j} \mathrm{X} ^ {r ^ {j}}.
\]

\begin{enumerate}
\def\labelenumi{\alph{enumi})}
\item
  证明：若 P、Q 是 K{[}X{]} 中的两个多项式，且 R = PQ，则我们有 \(\tilde{\mathbf{R}}(\mathbf{X}) = \tilde{\mathbf{P}}(\tilde{\mathbf{Q}}(\mathbf{X}))\) 。由此推出：要使 K{[}X{]} 中的两个多项式 F、G 满足 F 整除 G，必须而且只需 \(\hat{F}\) 整除 \(\hat{G}\) 。（要看出该条件是充分的，考虑 G 除以 F 的欧几里得除法。）
\item
  设 F 是 K{[}X{]} 中的不可约多项式，并设 G 是 K{[}X{]} 中的任意多项式。如果存在整数 \(d \geqslant 1\) 使得 \(\hat{\mathrm{F}}(\mathbf{X}^{d})\) 与 \(\hat{\mathrm{G}}(\mathbf{X}^{d})\) 有公共根，证明 F 整除 G。（若 H 是 \(\hat{\mathrm{F}}(\mathbf{X}^{d})\) 与 \(\hat{\mathrm{G}}(\mathbf{X}^{d})\) 的最大公因子，考虑多项式 \(P \in K[X]\) （H 整除 \(\hat{\mathrm{P}}(\mathbf{X}^{d})\) ）所成的集合。）
\item
  每个多项式 \(\mathbf{F} \neq 0\) （属于 \(\mathbf{K}[\mathbf{X}]\) ）都整除某个形如 \(\mathbf{X}^{\mathbf{N}} - 1\) 的多项式。证明：若 \(r - 1 = de\) 且 \(\alpha\) 是 \(\hat{\mathbf{F}}(\mathbf{X}^{d})\) 的一个根，则 \(\alpha\) 属于 \(\mathbf{K}\) 的次数为 \(n\mathbf{N}\) 的扩张，并由此推出 \(\hat{\mathbf{F}}(\mathbf{X}^d)\) 在 \(\mathbf{K}[\mathbf{X}]\) 中的每个不可约因子的次数都整除 \(n\mathbf{N}\) （利用 \(a\) ））。若进一步 \(e\) 与 \(d\mathbf{N}\) 互素，且 \(m\) 是 \(\alpha\) 在 \(\mathbf{K}\) 的次数为 \(n\) 的扩张上的次数，证明 \(m\) 整除 \(\mathbf{N}\) （利用 \((r^{\mathbf{N}} - 1)/e \equiv d\mathbf{N} (\text{mod } e)\) 这一事实）。
\end{enumerate}

\(d\) ) 采用与 \(c\) ) 相同的记号，假设 \(\mathbf{F}\) 不可约，且 \(\mathbf{N}\) 是整数 \(h\) 中最小的一个，其中 \(\mathbf{F}\) 整除 \(\mathbf{X}^h - 1\) 。若 \(e\) 与 \(d\mathbf{N}\) 互素，且 \(n\) 的每个素因子都整除 \(\mathbf{N}\) ，证明 \(\hat{\mathbf{F}}(\mathbf{X}^d)\) 的每个不可约因子的次数都是 \(n\mathbf{N}\) 。（注意，根据 \(b\) )，此时我们有 \(m = \mathbf{N}\) 。）

\begin{enumerate}
\def\labelenumi{\alph{enumi})}
\setcounter{enumi}{4}
\tightlist
\item
  特别地，由 \(d\) ) 推出：若 \(\mathrm{P}(\mathbf{X}) = \sum_{j} a_j \mathbf{X}^j\) 是
\end{enumerate}

K{[}X{]} 中的不可约多项式，且 N 是使 \(P(X)\) 整除 \(X^{h}-1\) 的最小整数 h ，则 \(\sum a_{j}X^{q^{j}-1}\) 的每个不可约因子的次数都是 N。

\begin{enumerate}
\def\labelenumi{\arabic{enumi})}
\setcounter{enumi}{9}
\tightlist
\item
  设 \(t, n, k\) 为三个 \(\geqslant 1\) 的整数，使得 \(t \leqslant n\) 且 \(n\) 整除 \(2^k - 1\) 。设 \(S\) 为满足下列条件的最小整数集合：\((1, t) \subset S \subset (1, n)\) ，且条件 \(\langle i \in S, 2i \equiv j (\text{mod } n) \text{ 且 } 1 \leqslant j \leqslant n \rangle\) 蕴含 \(j \in S\) 。令 \(m = \text{Card } S\) 。设 \(\xi\) 是一个阶为 \(n\) 的元素，含于 \(F_2^k\) ，并令 \(P(X) = \prod_{j \in S} (X - \xi^j)\) 。
\end{enumerate}

\begin{enumerate}
\def\labelenumi{\alph{enumi})}
\item
  证明 \(\mathbf{P}(\mathbf{X})\in \mathbf{F}_2[\mathbf{X}]\) 。
\item
  证明：若 \(\mathbf{R}(\mathbf{X})\in \mathbf{F}_2[\mathbf{X}]\) 的次数 \(< n\) 且是 \(\mathrm{P}(\mathbf{X})\) 的非零倍元，则其非零系数严格多于 \(t\) 个。
\item
  由此构作一个映射 \(\varphi: \mathbf{F}_2^n \to \mathbf{F}_2^m\) ，使得当 \(u, v \in \mathbf{F}_2^n\) 互异且至多有 \(t\) 个不同坐标时，有 \(\varphi(u) \equiv \varphi(v)\) 。
\item
  当 t = 6、n = 15、k = 4 时，计算 S、m 和 P。
\end{enumerate}

\begin{enumerate}
\def\labelenumi{\arabic{enumi})}
\setcounter{enumi}{10}
\tightlist
\item
  设 \(f\) 是 \(\mathbf{Z}[\mathbf{X}]\) 中的首一多项式，在 \(\mathbf{Q}[\mathbf{X}]\) 中不可约，并设 \(p\) 是一个素数，使得 \(\vec{f}\) —— \(f\) 在 \(\mathbf{F}_p[\mathbf{X}]\) 中的典范像 —— 是一个无重根的多项式。若 \(\Gamma\) （相应地 \(\vec{\Gamma}\) ）是多项式 \(f\) （相应地 \(\vec{f}\) ）在 \(\mathbf{Q}\) （相应地 \(\mathbf{F}_p\) ）上的伽罗瓦群，试给出从 \(\Gamma\) 的一个子群到 \(\vec{\Gamma}\) 之上的一个同构。由此推出：若 \(\vec{f} = g_1g_2\ldots g_r\) ，其中各 \(g_j\) 在 \(\mathbf{F}_p[\mathbf{X}]\) 中不可约，且 \(g_j\) 的次数为 \(n_j\) （从而 \(f\) 的次数为 \(n = n_1 + n_2 + \dots + n_r\) ），则群 \(\Gamma\) 作为 \(\mathfrak{S}_n\) 的子群，包含一个由支集两两不相交的循环组成的乘积 \(\sigma_1\sigma_2\ldots\sigma_r\) ，其中 \(\sigma_j\) 的阶为 \(n_j\) ，\(1\leqslant j\leqslant r\) （参见 I，第 62 页）。
\end{enumerate}

推广到 f 不是首一多项式、但其首项系数不被 p 整除的情形。

\begin{enumerate}
\def\labelenumi{\arabic{enumi})}
\setcounter{enumi}{11}
\item
  设 \(f\) 是 \(\mathbf{Z}[\mathbf{X}]\) 中次数为 \(n\) 的多项式，并设 \(p_1, p_2, p_3\) 是三个不整除 \(f\) 的首项系数的互异素数；设 \(f_1, f_2, f_3\) 分别是 \(f\) 在 \(\mathbf{F}_{p_1}[\mathbf{X}], \mathbf{F}_{p_2}[\mathbf{X}], \mathbf{F}_{p_3}[\mathbf{X}]\) 中的典范像。假设 \(f_1\) 是一个线性因子与一个次数为 \(n - 1\) 的不可约因子之积，\(f_2\) 是一个二次不可约因子与若干奇数次不可约因子之积，最后 \(f_3\) 不可约。证明：在这些条件下，\(f\) 在 \(\mathbf{Q}[\mathbf{X}]\) 中不可约，且多项式 \(f\) 在 \(\mathbf{Q}\) 上的伽罗瓦群同构于对称群 \(\mathfrak{S}_n\) （«戴德金准则»；同时利用习题 11 以及 I，第 63 页，命题 9）。
\item
  设 \(p_1, p_2, \ldots, p_m\) 为互不相同的奇素数，且 \(P = p_1p_2 \ldots p_m\) （ \(m \geqslant 3\) ）。
\end{enumerate}

\begin{enumerate}
\def\labelenumi{\alph{enumi})}
\tightlist
\item
  给定整数 \(n \geqslant 1\) ，设 \(E\) 为 \(\mathbf{Z}[\mathbf{X}]\) 中次数 \(\leqslant n\) 、系数落在区间 \([0, P[\) （属于 \(\mathbf{N}\) ）内的多项式组成的集合；于是我们有 \(\text{Card}(E) = P^{n+1}\) 。设 \(k\) 是满足 \(0 < k < 1/3n\) 的数；证明：在多项式 \(f \in E\) 之中，至少有 \(k^m\mathbf{P}^{n + 1}\) 个使得对每个指标 \(j\) （满足 \(1 \leqslant j \leqslant m\) ），典范像 \(f_j\) —— \(f\) 在 \(\mathbf{F}_{p_j}[\mathbf{X}]\) 中的 —— 的次数为 \(n\) 且不可约（利用习题 1，b））。再假设 \(k < \frac{1}{18(n - 2)}\) 且 \(n > 2\) ；以同样的方式论证并利用戴德金准则（习题 12）
\end{enumerate}

证明存在一个由多项式 \(f \in \mathbf{E}\) 组成的集合，其基数至少等于 \((1 - 3(1 - k)^m) \mathbf{P}^{n+1}\) ，这些多项式次数为 \(n\) 、不可约，且其在 \(\mathbf{Q}\) 上的伽罗瓦群同构于 \(\mathfrak{S}_n\) 。

\begin{enumerate}
\def\labelenumi{\alph{enumi})}
\setcounter{enumi}{1}
\tightlist
\item
  对每个整数 N，设 \(L_{n,N}\) 为 Z{[}X{]} 中次数 \(\leqslant n\) 、系数属于区间 \((-N, N)\) 的多项式组成的集合；于是我们有 \(\text{Card}(L_{n,N}) = (2N + 1)^{n+1}\) 。固定整数 n，证明：对每个 \(\varepsilon\) （满足 \(0 < \varepsilon < 1\) ），存在整数 \(N_0\) ，使得当 \(N \geqslant N_0\) 时，多项式 \(f \in L_{n,N}\) 中至少有 \((1 - \varepsilon)(2N + 1)^{n+1}\) 个是次数为 n 的不可约多项式，且其在 Q 上的伽罗瓦群同构于 \(S_n\) （利用 a））。
\end{enumerate}

* 14) 设 K 是一个有限域，交换与否均可；设 Z 是它的中心，q 是 Z 的元素个数，n 是秩 {[}K:Z{]} 。

\begin{enumerate}
\def\labelenumi{\alph{enumi})}
\item
  若 E 是 K 的包含 Z 的子域，证明 \([E:Z]\) 整除 n。
\item
  设 \(x \in K\) 是不属于 Z 的元素；证明 x 的不同共轭元 \(yxy^{-1}\) 在 \(K^{*}\) 中的个数具有形式 \((q^{n}-1)/(q^{d}-1)\) ，其中 d 整除 n 且异于 n（在 K 中考虑与 x 可交换的元素组成的集合，并利用 a））。
\item
  由 \(b\) ) 推出 \(q - 1\) 被整数 \(\Phi_n(q)\) 整除（把群 \(\mathbf{K}^*\) 划分成共轭类，并利用 \(V\) ，第 82 页的关系式 (6)）。
\item
  证明：若 \(n > 1\) ，则 \(\Phi_n(q) > (q - 1)^{\varphi(n)}\) （把 \(\Phi_n(X)\) 在复数域 \(C\) 中分解）。由此推断必有 \(K = Z\) ，换言之，\(K\) 必然是交换的（韦德伯恩定理）。*
\end{enumerate}

